\section{Discussion}

\subsection{Differences Between Bac and Vir Networks}

The observed differences in community structure between the Bac and Vir networks
can primarily be attributed to their underlying structural organization.
The Bac network is characterized by a higher edge density and stronger overall connectivity,
suggesting the presence of tightly interconnected protein communities.
In contrast, the Vir network exhibits a sparser structure with a larger 
proportion of weakly connected nodes.

\subsection{Which Method Is More Stable}

Stability in community detection refers to the robustness 
of identified partitions under methodological variation.
The high agreement between Louvain and Leiden in both networks, as indicated by the ARI values,
suggests that both methods produce stable community assignments.
This reinforces confidence in the reliability of the detected communities.

\subsection{Resolution Limit Problem}

An important limitation of modularity-based community detection methods, such as Louvain and Leiden,
is the resolution limit problem. Modularity optimization may fail to 
detect smaller but well-defined communities if merging them with larger groups
results in a higher overall modularity score.

This limitation is particularly relevant for dense networks such as Bac, where small modules
may be overlooked in favor of larger communities, leading to a
less accurate representation of the true community structure.

Although the Leiden algorithm is designed to produce better-connected communities, 
it is not immune to the resolution limit problem.
Therefore, the detected partitions should be interpreted as optimal with respect to the 
chosen resolution parameter rather than as uniquely defined ground-truth communities.